\documentclass[a4paper,12pt]{letter}

\usepackage{amsmath}
\usepackage{amssymb}
\usepackage{nopageno}
\usepackage{amsmath}
\usepackage{enumitem}
% \usepackage{mathtools}
%\usepackage{eqnarray,amsmath}
\usepackage[a4paper, total={7.5in, 11.25in}]{geometry}
%\usepackage[margin=2.5cm]{geometry}
\usepackage[utf8]{inputenc}
\usepackage[T1]{fontenc}
\usepackage{lmodern}
\usepackage{amsmath, amssymb}
\usepackage{setspace}
\setstretch{1.2}


\begin{document}

\begin{center}
  (KTXFI2EBNF) Physics II. Homework No.~4 Solutions\\ \smallskip
 December~18,~2025
 \end{center}


\begin{enumerate}

\item By how much does the emitted current change if the tungsten cathode temperature is increased from $T = 2250\,\text{K}$ to $2400\,\text{K}$?  The work function for tungsten is $W_{\text{ki}} = 5.0\,\text{eV}$.  The Richardson coefficient is $6.05 \cdot 10^{5}\,\text{A}/\text{m}^2\text{K}^2$.

  \medskip

The emission current density is given by the Richardson--Dushman equation $$ j = A T^2 \exp\!\left(-\frac{W_{ki}}{k T}\right).$$

The ratio of the emitted current densities is $$ \frac{j_2}{j_1} = \frac{T_2^2}{T_1^2} \exp\!\left[-\frac{W_{ki}}{k}\left(\frac{1}{T_2} - \frac{1}{T_1}\right)\right].$$

Substituting the numerical values, $$ \frac{j_2}{j_1}=\frac{2400^2}{2250^2}\exp\!\left[-\frac{5.0 \cdot 1.6 \cdot 10^{-19}}{1.38 \cdot 10^{-23}}\left(\frac{1}{2400} - \frac{1}{2250}\right)\right].$$

The emitted current increases by approximately a factor of $$\frac{j_2}{j_1} \approx 8.$$

% -----------------------------------------------------------------------------------------------------

\item What retarding voltage reduces the emission current density to one fifteenth for a cathode at $T = 2300\,\text{K}$?  
    The electron charge is $q = -1.6 \cdot 10^{-19}\,\text{As}$.

    \medskip

The current density in the presence of a retarding voltage $U$ is
$$
j = j_0 \exp\!\left(\frac{q U}{k T}\right).
$$

The condition
$$
\frac{j}{j_0} = \frac{1}{15}
$$
leads to
$$
\exp\!\left(\frac{q U}{k T}\right) = \frac{1}{15}.
$$
Taking the natural logarithm,
$$
\frac{q U}{k T} = \ln\!\left(\frac{1}{15}\right).
$$

Solving for the retarding voltage,
$$
U = \frac{k T}{q} \ln\!\left(\frac{1}{15}\right).
$$

Substitution yields
$$
U =
\frac{1.38 \cdot 10^{-23} \cdot 2300}
{1.6 \cdot 10^{-19}}
\ln\!\left(\frac{1}{15}\right),
$$
which gives
$$
U \approx -0.55\ \mathrm{V}.
$$

\bigskip

% -----------------------------------------------------------------------------------------------------

\pagebreak
  \item What is the free-electron drift velocity in a copper conductor of radius $R = 2.5\,\text{mm}$ if the current is $I = 15\,\text{A}$?  
  The electron density in copper is $n_{\text{e}} = 3.3 \cdot 10^{28}\,\text{m}^{-3}$.

\medskip

The current is related to the drift velocity by $$I = n_e q v_d A, $$ where the cross-sectional area is $$A = \pi R^2.$$

Solving for the drift velocity, $$v_d = \frac{I}{n_e q \pi R^2}.$$

Substituting the given values,$$v_d =\frac{15}{3.3 \cdot 10^{28} \cdot 1.6 \cdot 10^{-19}\cdot \pi \cdot (2.5 \cdot 10^{-3})^2}.$$

The drift velocity is $$v_d \approx 1.4 \cdot 10^{-4}\ \mathrm{m/s}.$$

\bigskip

\item We examine an n-type semiconductor under the conditions shown in the figure:  $B_{y} = 0.5\,\text{T}$, $U_{\text{H}} = 12\,\text{mV}$, $I = 25\,\text{mA}$.  The magnetic induction is directed along the $y$–axis.


  \begin{enumerate}
  \item What is the electron concentration?

  \medskip

The Hall voltage is given by $$U_H = \frac{I B_y}{n q d}, $$ where $d$ is the sample thickness.

Solving for the electron concentration, $$n = \frac{I B_y}{q U_H d}.$$

\medskip

% -----------------------------------------------------------------------------------------------------

\item What is the drift velocity of the electrons?

  \medskip

The drift velocity follows from the Hall field: $$ v_d = \frac{U_H}{B_y d}.$$

\medskip

% -----------------------------------------------------------------------------------------------------

\item What is the Hall constant?

  \medskip

The Hall constant is defined as $$ R_H = \frac{1}{n q}.$$

  
\end{enumerate}

\end{enumerate}



\bigskip

\rightline{Dr.~Gabor~FACSKO}
\rightline{\textit{facsko.gabor@uni-obuda.hu}}

\vfill

\end{document}
